\documentclass{article}

\usepackage{tabularx}
\usepackage{booktabs}

\title{Problem Statement and Goals\\\progname}

\author{\authname}

\date{}

\input{../Comments.text}
\input{../Common.text}

\begin{document}

\maketitle

\begin{table}[hp]
\caption{Revision History} \label{TblRevisionHistory}
\begin{tabularx}{\textwidth}{llX}
\toprule
\textbf{Date} & \textbf{Developer(s)} & \textbf{Change}\\
\midrule
Date1 & Name(s) & Description of changes\\
Date2 & Name(s) & Description of changes\\
... & ... & ...\\
\bottomrule
\end{tabularx}
\end{table}

\section{Problem Statement}

\wss{You should check your problem statement with the
\href{https://smiths.github.io/capTemplate/Checklists/ProbState-Checklist.pdf}
{problem statement checklist}}. 

\wss{You can change the section headings, as long as you include the required
information.}

\subsection{Problem}

\wss{What problem are you trying to solve?  Why is it important? Do not talk
about how you will solve the problem, only what the problem is.  AI is rarely
part of the problem.  The problem is ``what'' you want to do, not ``how'' you
want to do it.  The problem statement should be written in a way that is
understandable to a non-technical audience.}

\subsection{Inputs and Outputs}

\wss{Characterize the problem in terms of ``high level'' inputs and outputs.  
Use abstraction so that you can avoid details.}

\subsection{Stakeholders}

\subsection{Environment}

\wss{Hardware and Software Environment for the users.  The developer environment
is summarized as part of the developer plan.}

\section{Goals}

\section{Stretch Goals}

\wss{Goals you hope to get to, but not necessary.  They are also helpful if the
scope of the original project needs to be expanded.}

\section{Custom Documents (CDs)}

\wss{For CAS 741: State whether the project is a research project. This
designation, with the approval (or request) of the instructor, can be modified
over the course of the term.}

\wss{For SE Capstone: List your custom documents (CDs).  Potential CDs include
usability testing, code walkthroughs, user documentation, formal proof,
GenderMag personas, Design Thinking, etc.  (The full list is on the course
outline and in Lecture 02.) Normally the number of CDs will be two (one per
term).  Approval of the CDs will be part of the discussion with the instructor
for approving the project. The CDs, with the approval (or request) of the
instructor, can be modified over the course of the term.}

\newpage{}

\section*{Appendix --- Reflection}

\input{../Reflection.text}

\begin{enumerate}
    \item What went well while writing this deliverable?
    \item What pain points did you experience during this deliverable, and how
    did you resolve them?
    \item How did you and your team adjust the scope of your goals to ensure
    they are suitable for a Capstone project (not overly ambitious but also of
    appropriate complexity for a senior design project)?
\end{enumerate}  
\noindent\rule{\linewidth}{0.4pt}
\begin{enumerate}
    \item This deliverable went smoothly, and a key reason was because our team
has done a lot of discussion about what we wanted our Capstone experience to be
like prior to the start of the term. Many of our prior conversations involved
someone bringing up a potential idea for the project, followed by the rest of us
questioning its details to try to understand it completely. That's why we've found
that, when it came time to start putting down documentation for these deliverables,
the answers came quickly because we've already discussed many of these topics.

    \item Despite our regular discussions on project decisions, we still found
that there were aspects of the project we hadn't yet considered but disagreed on.
For example, when it came to choosing our desired Git workflow, some members wanted to
use a feature branch workflow while others wanted to use a forking workflow. Both
sides felt pretty strongly about their choice, so we debated at length on which workflow
was better. We eventually decided upon a middle ground by using both workflows after
an impassioned discussion, and we made sure to resolve things nicely by reminding each
other of a few things: 1, nothing is personal when we're working through project stuff,
and 2, working through \& talking about disagreements is a better option than brushing
things under the rug.

    \item Scope was an important aspect for our team to consider because for many
of us, this was our first foray into working with technologies surrounding LLMs.
Due to this, we initially didn't have a great idea on what would be a feasible project
for Capstone. To educate ourselves, we looked into related projects and research papers
to get a grasp on the main concepts behind our distributed LLM idea. Once we understood
the core concepts behind the technology and learned about the new ideas that related
projects were bringing to the table, we broke down the project into manageable sections:
the core functionality needed for a minimum viable product, and extensions to the core
that make the project more interesting. This ensures that the scope of the project is not overly
ambitious because the core is minimal, but it also ensures appropriate complexity because
we have many potential extensions that could add more functionality, compatibility, or
usability to the final project.
\end{enumerate}  

\end{document}